
\emph{est de type fini.}

En effet, on peut supposer \'evidemment @@M0@@ r\'eduit et irr\'eductible d'apr\`es 3.4 (i). Alors il existe un ouvert non-vide @@M1@@ de @@M2@@ tel que @@M3@@ et @@M4@@ soient repr\'esentables et que la restriction de @@M5@@ soit quasi-affine (XII 1.1), donc de type fini par 3.2 (i). Munissant @@M6@@ de la structure r\'eduite induite, on peut supposer que le morphisme @@M7@@ induit par @@M8@@ est de type fini par r\'ecurrence noeth\'erienne\,; d'o\`u l'assertion, gr\^ace \`a 3.4 (i) appliqu\'e \`a @@M9@@.

\textbf{Th\'eor\`eme 3.6.} \emph{Soient @@M10@@ un sch\'ema noeth\'erien, @@M11@@ un @@M12@@-sch\'ema propre et @@M13@@ le foncteur de Picard. Alors pour tout entier @@M14@@, l'homomorphisme de puissance @@M15@@-i\`eme}
@@M16@@
\emph{est un morphisme de type fini.}

En effet, soit @@M17@@ une pr\'esentation de Chow~: @@M18@@ est projectif sur @@M19@@ et @@M20@@ est surjectif. Si l'on admet le th\'eor\`eme pour @@M21@@, alors en vertu de 3.5 et 3.4 (ii), le th\'eor\`eme pour @@M22@@ en r\'esulte. On peut donc supposer @@M23@@ projective.

Il existe un ouvert non-vide @@M24@@ de @@M25@@, un sch\'ema int\`egre @@M26@@ et un morphisme fini et surjectif @@M27@@, tels que toute composante irr\'eductible @@M28@@ de @@M29@@, munie de la structure r\'eduite induite, soit int\`egre sur @@M30@@ (XII 1.3). Munissant @@M31@@ de la structure r\'eduite induite, on peut admettre le th\'eor\`eme pour @@M32@@ par r\'ecurrence noeth\'erienne. Si l'on l'admet pour tous les @@M33@@, il en r\'esulte pour @@M34@@ par le
